\section{Methodology}
\label{sec:methodology}

\subsection{Datasets}
\label{subsec:method_problem_data}

We used three retinal fundus datasets and one microscopy dataset, which differ in image size and in the number of training images. DRIVE is a standard benchmark of retinal vessel segmentation~\cite{moccia2018blood}, with few, small images; FIVES has many images at high resolution; HRF adds a second high-resolution dataset; and VessMap is a second imaging domain, with few, small and heterogeneous images. The DRIVE and VessMap images are smaller than the window of \medsam{}, which sees them only after upscaling. DRIVE~\cite{Staal2004} has 40 images of $565 \times 584$ pixels~\cite{galdran2020littlewnetcouldstateoftheart} with an official split of 20 training and 20 test images; we held out four of the training images (images 37 to 40) for validation and trained on the other 16. HRF~\cite{budai2013robust} has 45 images of $3504 \times 2336$ pixels, and FIVES~\cite{jin2022fives} has 800 images of $2048 \times 2048$ pixels with an official split of 600 training and 200 test images. VessMap~\cite{viana2025new} has confocal microscopy images of the mouse brain cortex vasculature ($256 \times 256$ pixels), with low contrast; it is a check in a second imaging domain, and the models were trained and tuned on it like on the other datasets. HRF and VessMap have no official split, so we split each at random, with a fixed seed, into 70\% training and 30\% test images: 31 training and 14 test images for HRF, and 70 and 30 for VessMap. On HRF, FIVES and VessMap no separate validation set was held out, and checkpoint selection, early stopping and the hyperparameter search used the test images. This protocol was the same for every model. Appendix~\ref{app:training} describes a correction of the border pixels of some FIVES masks.

\subsection{Models}
\label{subsec:method_models}
\label{subsec:method_matrix}

\begin{table}[tb]
\caption{The five architectures and their parameters. The size of \convunext{} depends on the dataset (Appendix~\ref{app:training}).}
\label{tab:model_matrix}
\begin{ruledtabular}
\begin{tabular}{lllr}
Model & Type & Encoder & Params (M) \\
\midrule
\medsam{} & Plain ViT & ViT-B & 93.7 \\
\litemedsam{} & Hybrid & TinyViT & 9.8 \\
\convnext{} & CNN & ConvNeXt-L & 211.4 \\
\convunext{} & CNN & U-Net & 3.5 / 13.8 \\
\lwnet{} & CNN & W-Net & 0.07 \\
\end{tabular}
\end{ruledtabular}
\end{table}

Table~\ref{tab:model_matrix} lists the architectures. \medsam{}~\cite{ma2024segment} is the plain ViT, a single-scale ViT without convolutional stages: the ViT-B encoder of SAM, with 16~px patches and a $1024 \times 1024$~px input window, pre-trained on over one million medical image--mask pairs. Its fundus images come from optic disc and cup datasets, and none of our four datasets is among its training data~\cite{ma2024segment}. \litemedsam{}~\cite{ma2024efficient} is the hybrid: its TinyViT encoder~\cite{wu2022tinyvit} combines convolutional stages (mobile inverted bottleneck convolutions) with windowed attention, on 256~px windows; it was distilled and fine-tuned on the same data as \medsam{}~\cite{ma2024efficient}, so neither model was trained on our datasets. Both transformer-based models keep SAM's prompt encoder and mask decoder~\cite{Kirillov_2023_ICCV}. In training and inference, the prompt is a fixed box covering the whole input window, and the prompt encoder is not trained. The decoder upsamples the $64 \times 64$ grid of encoder tokens four times, so \medsam{} predicts at a quarter of its input resolution and its logits are upsampled bilinearly to the window, while \litemedsam{}, whose encoder reduces its 256~px window to the same grid, predicts at its input resolution. The three CNNs span more than three orders of magnitude in size: \convnext{}, a ConvNeXt-Large encoder~\cite{liu2022convnet} with a transposed-convolution decoder; \convunext{}, a U-Net built from ConvNeXt-style blocks with attention gates~\cite{han2022convunext}; and \lwnet{}, our implementation of Little W-Net~\cite{galdran2020littlewnetcouldstateoftheart}.

\subsection{Training protocol}
\label{subsec:method_training}

We trained the three CNNs from random initialization. The two transformer-based models were trained in two ways: fine-tuned from their published pre-trained weights~\cite{ma2024segment, ma2024efficient}, with all weights updated except those of the prompt encoder, and trained from random initialization. We also trained a compact version of each transformer-based model from random initialization, with a 256~px input and an architecture size chosen by the hyperparameter search on each dataset (Table~\ref{tab:compact_models}).

For \medsam{}, the compact version also changes how the image is split into tokens. \medsam{} divides its 1024~px window into 16~px patches, a $64 \times 64$ grid of tokens; compact \medsam{} works on a 256~px window, and the search selected 4~px patches on DRIVE, HRF and VessMap ($64 \times 64$ tokens) and 8~px patches on FIVES ($32 \times 32$ tokens). A token of compact \medsam{} therefore covers a 4 or 8~px square of the native image, against 16~px for \medsam{} on HRF and FIVES and about 9 and 4~px on DRIVE and VessMap, which \medsam{} sees upscaled. Compact \litemedsam{} keeps the 256~px input of \litemedsam{} and differs from it only in size.

Comparisons across architectures pair the CNNs, trained from scratch, with the fine-tuned transformer-based models, so this comparison basis mixes initializations. Each configuration was trained once, with one random seed, on each dataset.

Each training sample was a random crop of the model's input size, 256~px, or 1024~px for \medsam{}, taken after random flips and rotations. Images smaller than the input, which is the case for DRIVE and VessMap with \medsam{}, were upscaled to it, so \medsam{} saw DRIVE at about 1.75 times and VessMap at 4 times their native scale. All models were trained with the soft Dice loss~\cite{milletari2016vnet} and AdamW~\cite{loshchilov2019decoupled}, in PyTorch~\cite{paszke2019pytorch} (version 2.10.0) with PyTorch Lightning~\cite{falcon2026lightning} (version 2.6.1). For each configuration and dataset, a search of 100 trials chose the learning rate, the weight decay, the batch size, one architecture parameter of each CNN and the architecture of the compact models (Table~\ref{tab:search_space}), and we kept the trial with the highest validation Dice. Training stopped after 50 epochs without improvement of the validation Dice, or at 1000 epochs. Appendix~\ref{app:training} gives the full training details.

\subsection{Evaluation protocol and metrics}
\label{subsec:method_metrics}

We evaluated every model on the full test images at native resolution with the sliding-window inference of MONAI~\cite{cardoso2022monai} (version 1.5.2): windows of the model's input size with 25\% overlap, combined with Gaussian weighting, and images smaller than a window padded to it. The sigmoid output was binarized at 0.5. For \medsam{} on DRIVE and VessMap, which saw upscaled images during training, we upscaled each test image to the $1024 \times 1024$~px input, ran one window, resized the logits back to the native resolution with bilinear interpolation and scored them against the native ground truth. We call this the upscaled protocol. It evaluates \medsam{} at its training scale, and we adopt it for every \medsam{} result. Table~\ref{tab:padding} compares it with the padded protocol, which feeds the native image padded to the window at the same cost. HRF and FIVES images are larger than the window, so the two protocols coincide on these datasets.

All metrics were computed for each test image from the binarized prediction $\hat{Y}$ and the ground truth $Y$. Dice is $2|Y \cap \hat{Y}| / (|Y| + |\hat{Y}|)$ over the full image. The centerline metrics use the skeleton $S(\cdot)$ of a mask, from the \texttt{skeletonize} function of scikit-image~\cite{vanderwalt2014scikit} (version 0.26.0; Zhang--Suen thinning~\cite{zhang1984fast}), and the radius of a skeleton pixel, from the Euclidean distance transform of the same mask, computed with SciPy~\cite{virtanen2020scipy}. The centerline recall $R_{\mathrm{cl}} = |S(Y) \cap \hat{Y}| / |S(Y)|$ is the share of ground-truth centerline pixels inside the predicted mask, and the centerline precision $P_{\mathrm{cl}} = |S(\hat{Y}) \cap Y| / |S(\hat{Y})|$ the share of predicted centerline pixels inside the ground-truth mask, with no distance tolerance. To measure recall as a function of vessel caliber, we grouped the ground-truth centerline pixels by radius into the bins $[0.5, 1.5)$, $[1.5, 2.5)$, $[2.5, 3.5)$, $[3.5, 4.5)$ and $[4.5, \infty)$~px. We call a vessel thin when its radius is below 1.5~px. The centerline recall of thin vessels is the recall of the first bin,
\begin{equation}
    \thinrecall = \frac{|S_{\mathrm{thin}}(Y) \cap \hat{Y}|}{|S_{\mathrm{thin}}(Y)|},
    \label{eq:thin_recall}
\end{equation}
where $S_{\mathrm{thin}}(Y)$ is the set of ground-truth centerline pixels with radius below 1.5~px. The centerline precision of thin vessels, $P_{\mathrm{thin}} = |S_{\mathrm{thin}}(\hat{Y}) \cap Y| / |S_{\mathrm{thin}}(\hat{Y})|$, takes the predicted centerline pixels whose radius in the prediction is below 1.5~px. The two thin-vessel metrics are therefore computed on different sets of pixels: $\thinrecall$ on the ground-truth thin centerline, $P_{\mathrm{thin}}$ on the predicted one. A value with an empty denominator is undefined and is left out of the mean.

% Generated by paper/figures_tables/tab01_global_baseline.py -- do not edit by hand.
% Source: results/per_image_metrics.csv (MedSAM DRIVE/VessMap upscaled protocol); tests from
% outputs/statistics/all_tests.csv.
\begin{table*}[tbp]
\caption{Overall segmentation performance. Bold: best; underlined: not significantly different from the best (per row).}
\label{tab:global_baseline}
\begin{ruledtabular}
\begin{tabular}{llccccc}
 & & \multicolumn{3}{c}{CNN} & \multicolumn{2}{c}{Fine-tuned} \\
\cmidrule(lr){3-5} \cmidrule(lr){6-7}
Dataset & Metric & \lwnet{} & \convunext{} & \convnext{} & \litemedsam{} & \medsam{} \\
\midrule
\multirow{3}{*}{\shortstack{DRIVE\\($n=20$)}} & Dice & $81.1 \pm 2.4$ & $77.6 \pm 15.1$ & $77.5 \pm 2.5$ & {\boldmath$82.3 \pm 1.9$} & $81.3 \pm 1.3$ \\
 & $R_{\mathrm{cl}}$ & $74.9 \pm 7.2$ & $70.1 \pm 16.7$ & $66.9 \pm 7.5$ & {\boldmath$79.4 \pm 7.1$} & $75.5 \pm 6.9$ \\
 & $P_{\mathrm{cl}}$ & \underline{$89.3 \pm 2.3$} & {\boldmath$90.5 \pm 4.4$} & \underline{$89.5 \pm 3.1$} & $87.6 \pm 2.5$ & $87.3 \pm 2.9$ \\
\midrule
\multirow{3}{*}{\shortstack{HRF\\($n=14$)}} & Dice & $76.9 \pm 6.9$ & $79.7 \pm 5.1$ & $76.9 \pm 6.2$ & \underline{$80.6 \pm 5.9$} & {\boldmath$81.7 \pm 4.5$} \\
 & $R_{\mathrm{cl}}$ & $72.0 \pm 12.0$ & $72.8 \pm 9.5$ & $73.8 \pm 7.6$ & {\boldmath$82.2 \pm 7.6$} & $76.6 \pm 8.6$ \\
 & $P_{\mathrm{cl}}$ & \underline{$84.0 \pm 11.5$} & \underline{$88.5 \pm 7.6$} & $79.5 \pm 10.0$ & $83.9 \pm 8.0$ & {\boldmath$88.9 \pm 6.7$} \\
\midrule
\multirow{3}{*}{\shortstack{FIVES\\($n=200$)}} & Dice & $81.8 \pm 14.7$ & $89.3 \pm 9.2$ & $90.2 \pm 8.5$ & $88.1 \pm 9.8$ & {\boldmath$92.2 \pm 8.5$} \\
 & $R_{\mathrm{cl}}$ & $78.7 \pm 17.3$ & $87.4 \pm 11.1$ & $88.9 \pm 10.0$ & $87.6 \pm 10.9$ & {\boldmath$89.6 \pm 10.4$} \\
 & $P_{\mathrm{cl}}$ & $88.0 \pm 11.1$ & $92.1 \pm 7.4$ & $92.1 \pm 7.7$ & $89.4 \pm 10.6$ & {\boldmath$95.4 \pm 4.7$} \\
\midrule
\multirow{3}{*}{\shortstack{VessMap\\($n=30$)}} & Dice & \underline{$92.8 \pm 2.4$} & \underline{$92.9 \pm 2.1$} & $91.0 \pm 2.3$ & {\boldmath$93.0 \pm 2.0$} & \underline{$92.8 \pm 2.1$} \\
 & $R_{\mathrm{cl}}$ & $95.9 \pm 3.6$ & $95.4 \pm 4.2$ & $94.7 \pm 4.1$ & {\boldmath$96.9 \pm 3.1$} & \underline{$96.5 \pm 3.7$} \\
 & $P_{\mathrm{cl}}$ & \underline{$96.2 \pm 3.1$} & {\boldmath$97.0 \pm 2.0$} & \underline{$96.3 \pm 3.1$} & \underline{$96.8 \pm 2.8$} & $96.3 \pm 2.7$ \\
\end{tabular}
\end{ruledtabular}
\end{table*}


We counted the parameters of each trained checkpoint and its floating-point operations with fvcore~\cite{fvcore2022} (version 0.1.5) for one forward pass over one window, counting one multiply--add as one FLOP. The cost per image is the cost per window times the number of windows over a full test image. We did not measure latency or throughput. As a measure of training difficulty, we report the number of epochs the selected run needed to first reach 95\% of its best validation Dice.

Tables report the mean $\pm$ standard deviation of each metric over the test images, in \%, and figures report the mean. We compared every pair of configurations on each dataset and metric with a two-sided Wilcoxon signed-rank test~\cite{wilcoxon1945individual} on the per-image paired differences, with the Holm correction~\cite{holm1979simple} over the pairs of that dataset and metric. For Dice, $\thinrecall$ and $P_{\mathrm{thin}}$ these are the 36 pairs of the nine configurations; for $R_{\mathrm{cl}}$ and $P_{\mathrm{cl}}$, the 10 pairs of the five models of Table~\ref{tab:global_baseline}; and the two protocols of \medsam{} form one family of 12 tests (Table~\ref{tab:padding}). A difference is called significant when its Holm-adjusted p-value is below 0.05; Appendix~\ref{app:statistics} gives the computation and every p-value. The tests take the test images as the sample; with one training run per configuration, they do not include the variation between training runs. In the tables, bold marks the best value of a comparison (the highest mean, or the fewest epochs in Table~\ref{tab:medsam_vs_cnn}); where more than two values are compared, underline marks every value not significantly different from it; and $^{*}$ marks a significant difference $\Delta$, in percentage points.
